La minería de datos o data mining, hace referencia a todas las técnicas que permiten extraer conocimiento de un conjunto de datos para descubrir, si es que existen, relaciones, modelos, regularidades o patrones que subyacen en un conjunto de datos y que, a priori, desconocemos.
\begin{figure}[H]
\centering
\begin{tikzpicture}[
    % Estilo para los bloques (Datos, Información, Conocimiento)
    box/.style={
        rectangle, 
        rounded corners=4pt, 
        draw=black, 
        fill=white, 
        fill opacity=0.95, 
        text opacity=1,
        minimum width=2.2cm,
        inner sep=8pt,
        font=\sffamily\bfseries\small,
        align=center
    },
    % Estilo para las flechas curvas de flujo
    arrow style/.style={
        -{Stealth[scale=1.3]},
        line width=1.5pt,
        draw=gray!70
    },
    % Estilo para las etiquetas de las flechas
    label style/.style={
        font=\sffamily\scriptsize\bfseries,
        align=center,
        text=black!70
    }
]

    % Triángulo de fondo (Cadena de valor)
    \fill[gray!30] (-0.5,-1.8) -- (-0.5,1.8) -- (6.5,0) -- cycle;

    % Bloques superpuestos
    \node[box, minimum height=3.2cm] (datos) at (0.6,0) {Datos};
    \node[box, minimum height=2.2cm] (info) at (2.8,0) {Información};
    \node[box, minimum height=1.2cm, minimum width=2.4cm] (conoc) at (5.2,0) {Conocimiento};

    % Flecha curva superior (Descubrimiento del valor)
    \draw[arrow style] (0.6, 1.9) to[bend left=30] (5.2, 0.9);
    \node[label style, above=0.15cm] at (2.9, 1.9) {DESCUBRIMIENTO DEL VALOR};

    % Flecha curva inferior (Aprovechamiento del valor)
    \draw[arrow style] (5.2, -0.9) to[bend left=30] (0.6, -1.9);
    \node[label style, below=0.15cm] at (2.9, -1.9) {APROVECHAMIENTO DEL VALOR};

\end{tikzpicture}
\caption*{\small \textbf{Figura 1.3.} Cadena de valor del dato.}
\end{figure}

Abarca técnicas de limpieza, transformación y análisis de datos con el fin de descubrir patrones y características específicas. Mientras que el aprendizaje automático engloba los algoritmos y modelos para que las máquinas aprendan y resuelvan tareas automáticamente.